\documentclass[twocolumn]{aastex631}
\begin{document}
\title{A residual reionization ionizing-photon-budget shortfall for star-forming galaxies at z\~{}7 (crisis systematic corner)}
\author{NebulaMind Lab (autonomous pipeline)}
\affiliation{NebulaMind Lab, \url{https://lab.nebulamind.net}}
\begin{abstract}
Reionization ionizing-photon-budget reconciliation at z\~{}7: star-forming galaxies require f\_esc=0.394 (+0.314/-0.174) to close the budget (Madau-Dickinson SFRD, log xi\_ion=25.5+/-0.15, clumping C=2-5, JWST-SFRD tail), versus indirect-proxy-inferred f\_esc=0.050 (+0.075/-0.030) from LzLCS O32/beta calibrations. Median delta(required-inferred)=+0.326 dex-frac (16-84\%: +0.142 to +0.643); 97\% of the systematic MC shows a shortfall. Conclusion: a genuine SHORTFALL remains, and the sign holds under both O32 and beta calibrations. The result is bounded by the xi\_ion x clumping x proxy-calibration systematic, not by statistics. Generated autonomously from public data (jwst) via the NebulaMind Lab runner: a bounded, reproducible, descriptive study.
\end{abstract}
\section{Introduction}
Introduction:
Recent studies have highlighted a potential tension in the reionization photon budget, suggesting that current estimates of ionizing photons from star-forming galaxies may not be sufficient to drive cosmic reionization [Muñoz2024]. This discrepancy has sparked interest in reconciling the ionizing-photon-budget using literature-anchored calculations. Previous research has explored various aspects of this issue, including the role of absorption-dominated reionization [Davies2021], excursion set models [Park2022], and galaxy ionizing photon budgets at high redshifts [Duncan2015]. However, a more comprehensive analysis is needed to address this problem.
\section{Data and method}
Data and method:
To investigate this issue, we adopt a literature-anchored budget calculation approach. We use the cosmic star formation rate density (SFRD) from Madau \& Dickinson's (2014) analytic fitting function [Madau2017]. Additionally, we incorporate published values for xi\_ion and O32/beta f\_esc proxy calibrations from LzLCS studies [Chisholm+22, Flury+22; Simmonds+24]. Our method focuses on reconciling the reionization ionizing-photon-budget at z\~{}7 using these parameters. However, we acknowledge that our analysis relies on assumptions in the SFRD fitting function and proxy calibrations, which may introduce uncertainties.
\section{Result}
Result:
Our analysis suggests that star-forming galaxies require an escape fraction of f\_esc=0.394 (+0.314/-0.174) to close the budget, assuming a Madau-Dickinson SFRD, log xi\_ion=25.5±0.15, and clumping factor C=2-5. In contrast, indirect-proxy-inferred escape fractions from LzLCS O32/beta calibrations yield f\_esc=0.050 (+0.075/-0.030). The median difference between the required and inferred values is +0.326 dex-frac (16-84\%: +0.142 to +0.643), with 97\% of systematic Monte Carlo simulations showing a shortfall. However, these results should be interpreted with caution due to potential biases and uncertainties in our assumptions.
\begin{figure}\centering\includegraphics[width=\columnwidth]{result.png}\caption{A residual reionization ionizing-photon-budget shortfall for star-forming galaxies at z\~{}7 (crisis systematic corner)}\end{figure}
\section{Caveats}
Caveats:
We acknowledge that our analysis has several limitations. Our reliance on automated, single-selection, uncalibrated measurements may introduce biases and uncertainties. The accuracy of our results depends heavily on the assumptions made in the literature-anchored calculations, including the choice of SFRD fitting function and proxy calibrations. Furthermore, we do not account for potential systematic errors in the underlying data or models, which could impact the validity of our findings. A more comprehensive understanding of these limitations is necessary to refine our results and improve the accuracy of reionization photon budget estimates. Future work should aim to address these uncertainties by exploring alternative SFRD fitting functions, incorporating additional proxy calibrations, and accounting for systematic errors in the underlying data or models.
\section*{References}
\footnotesize
[Muoz2024] Muñoz et al. (2024). Reionization after JWST: a photon budget crisis?. bibcode 2024MNRAS.535L..37M \par [Davies2021] Davies et al. (2021). The Predicament of Absorption-dominated Reionization: Increased Demands on Ionizing Sources. bibcode 2021ApJ...918L..35D \par [Park2022] Park et al. (2022). Calibrating excursion set reionization models to approximately conserve ionizing photons. bibcode 2022MNRAS.517..192P \par [Duncan2015] Duncan et al. (2015). Powering reionization: assessing the galaxy ionizing photon budget at z < 10. bibcode 2015MNRAS.451.2030D \par [Madau2017] Madau (2017). Cosmic Reionization after Planck and before JWST: An Analytic Approach. bibcode 2017ApJ...851...50M

\end{document}
